% Generated draft: requires mathematical and source review.
\documentclass[11pt]{article}
\usepackage[margin=25mm]{geometry}
\usepackage{fontspec}
\setmainfont{Latin Modern Roman}
\usepackage{amsmath,amssymb,mathtools,unicode-math}
\setmathfont{Latin Modern Math}
\usepackage{enumitem,needspace,xurl,hyperref}
\setlength{\parindent}{0pt}
\setlength{\parskip}{5pt}
\setlength{\emergencystretch}{3em}
\newcommand{\N}{\mathbb{N}}
\newcommand{\Z}{\mathbb{Z}}
\newcommand{\Q}{\mathbb{Q}}
\newcommand{\R}{\mathbb{R}}
\newcommand{\C}{\mathbb{C}}
\newcommand{\F}{\mathbb{F}}
\newcommand{\A}{\mathbb{A}}
\newcommand{\PP}{\mathbb{P}}
\newcommand{\E}{\mathbb{E}}
\newcommand{\eps}{\varepsilon}
\newcommand{\dd}{\,\mathrm{d}}
\begin{document}
\section*{UnsolvedMath ID: 6400003}
\textbf{Source problem code:} catalogue entry
\subsection{Definitions and mathematical statement}
A smooth four-manifold is a manifold of real dimension four equipped with a smooth structure. The source asks for a structure theory, or a conjectural decomposition, for such manifolds with a guiding role analogous to Thurston's Geometrization Conjecture in dimension three.

No specific decomposition theorem is asserted in the supplied record. A precise candidate theory would have to identify the class of four-manifolds covered, the allowable cutting or decomposition operations, the proposed building blocks, the reconstruction or gluing data, and the sense in which the resulting structure guides classification. The source does not fix any of these, or require uniqueness, finiteness, canonicity, a metric geometry, or an algorithm.

The dimension-three analogy is motivational: it does not impose a literal list of Thurston geometries or any particular cutting hypersurfaces in dimension four. The source also does not state compactness, boundary, orientability, or simple-connectivity assumptions, so none is added. Its target is the formulation of an organizing structure or conjecture, not proof of an unspecified universal classification theorem.
\subsection{Short English statement}
Develop a geometrization-like picture that can organize the study of smooth four-manifolds. The intended pieces, operations, scope, and classification criteria are part of the research program rather than given hypotheses.

The supplied summary says no complete analogous picture has emerged, while mentioning gauge-theoretic decomposition programs, elliptic surfaces, symplectic Lefschetz fibrations, and Donaldson's pencils/fibrations as partial organizing perspectives. Those symplectic tools are not asserted to apply to every smooth four-manifold. Its comments about active work in the 2020s and the program's open status are supplied assessments; provenance marks current status for review.
\subsection{Sources}
\begin{itemize}
\item Local UnsolvedMath registry, permanent ID 6400003, source code catalogue entry, ``A Geometrization Picture for Smooth Four-Manifolds,'' proposed by Simon Donaldson, 2008. Dataset curation and metadata: Ulam AI and UnsolvedMath Contributors, CC BY 4.0.
\item Simon Donaldson, \emph{Some Problems in Differential Geometry and Topology} (2008), explicit question/program 3 in the author PDF, \url{https://wwwf.imperial.ac.uk/~skdona/NONLINEARITYPROBLEMS.PDF}; normalized pdf-text extraction, accessed 2026-07-29.
\item The supplied summary invokes Thurston's geometrization as motivation and mentions gauge/symplectic building-block approaches without separate exact theorem references. Source-specific provenance marks current status and statement-reuse terms NEEDS\_REVIEW for the public author-hosted PDF, separately from the dataset licence.
\end{itemize}
\end{document}
